\chapter{数式データの扱い}\label{mexpr}

本章は，シンボリックな計算を扱うためのEgisonの機能を紹介する．

\section{数式データの内部表現}\label{mexpr-expr}

数式は組み込みデータ型としてEgison処理系内部でHaskellのデータ型を用いて直接実装されている．
本節は，数式データを表現するこの内部のデータ構造について説明する．

数式データは，分母と分子の両方に積和標準形の多項式をとる分数として表現されている．
そのため，たとえば，分母が多項式である分数同士の足し算の結果は，通分されて一つの分数にまとめられる．
\begin{lstlisting}[language=egison,numbers=none]
declare symbol x, y, a, b, c

a / (b + 1) + x / (y + 1)
-- (b x + a y + x + a) / (b y + y + b + 1)
\end{lstlisting}
多項式は，項のリストからなるデータとして表現されている．
そして，項は係数と，因子と冪指数の連想リストからなるデータとして表現されている．
たとえば，\lstinline{3 * x^2 * sqrt 2}という項は，係数として\lstinline{3}，連想リストとして\lstinline{[(x, 2), (sqrt 2, 1)]}をもつ．
冪指数には負の整数もとることができる．
つまり，Egisonの多項式は，ローラン多項式である．
分母がひとつの項だけからなる分数は，分数としては保持されず，負の冪指数を使って項のなかに取り込まれる．
\begin{lstlisting}[language=egison,numbers=none]
a / b + x / y -- x y^-1 + a b^-1
\end{lstlisting}

因子には，シンボルのほかに，\ref{cas-symbol2}節で紹介したシンボリックな関数適用，\ref{mexpr-backquote}節で紹介するクオートされた式，第\ref{funcsym}章で解説する関数シンボルをとることができる．
シンボルは，名前の文字列と添字のリストからなる．
添字については，テンソルについて解説する第\ref{tensor}章で説明する．

この内部構造は，Egisonライブラリで定義されている\lstinline{mathValue}マッチャーを使ったパターンマッチで分解できる．
たとえば，多項式を項のリストに分解できる．
\begin{lstlisting}[language=egison,numbers=none]
match (x - sqrt 2) as mathValue with
  | poly $ts -> ts
-- [x, - 'sqrt 2]
\end{lstlisting}
数式に対するパターンは，\ref{mexpr-pattern}節で詳しく紹介する．
また，同じ数式を複数の異なる内部構造（たとえば，平坦な多項式と入れ子の多項式）で表現し，型注釈でそれを選ぶこともできる．
これについては第\ref{casdetail}章で解説する．

数は複雑なデータ型である．
自然数までであれば代数的データ型としてエレガントに定義できる．
しかし，自然数の足し算の逆関数引き算を考えると，負の数という概念が得られる．
さらに足し算を繰り返すことによって得られる掛け算の逆を考えると，有理数という概念が得られる．
さらに掛け算を繰り返すことによって得られるべき乗という操作の逆を考えると，平方根や虚数，それらを含む代数的数の概念が得られる．

\section{バッククオート（\texttt{`}）による式展開の制御}\label{mexpr-backquote}

Egisonは数式を自動で積和標準形に展開すると\ref{cas-symbol}節で述べたが，この展開をバッククオート（\lstinline{`}）により制御することができる．
バッククオートに続く式は，ひとつのシンボルのように扱われる．
\begin{lstlisting}[language=egison,numbers=none,escapechar=']
`(x + 1)^2
--  `(x + 1)^2

`(x + 1) + `(x + 1)
--  2 * `(x + 1)

`(x + 1) + (x + 1)
-- x + `(x + 1) + 1
\end{lstlisting}

バッククオートが役に立つのは，主に多項式の割り算を含む計算のときである．
単変数の多項式については，\ref{cas-simplify}節で述べたユークリッドの互除法による約分が実装されているが，多変数の多項式の約分は実装されていない．
\begin{lstlisting}[language=egison,numbers=none,escapechar=']
(x + 1)^2 / (x + 1) -- x + 1
(x^2 - y^2) / (x - y) -- (- y^2 + x^2) / (- y + x)
\end{lstlisting}
バッククオートを使って因数分解された形を保っておけば，このような式も約分できる．
\begin{lstlisting}[language=egison,numbers=none,escapechar=']
`(x - y) * `(x + y) / `(x - y) -- `(y + x)
\end{lstlisting}

バッククオートによりクオートされた式は，Egison内部の数式データのなかで，シンボルと同列の因子として扱われる．
上記の\lstinline{`(x + 1)^2}が展開されずにひとつの因子の冪として保持されるのは，このためである．

\section{シングルクオート（\texttt{'}）による関数適用の制御}\label{mexpr-squote}

定義済みの関数の適用について，関数適用を評価せずに，シンボリックな結果を返したいことがある．
たとえば，\lstinline{sqrt}関数については\ref{cas-symbol2}節でみたように以下のように動作してほしい．
\begin{lstlisting}[language=egison,numbers=none]
declare symbol x, a

sqrt 4 -- 2
sqrt x -- 'sqrt x
\end{lstlisting}
このような動作を実現するために，Egisonには組み込み構文としてシングルクオート（\lstinline{'}）が用意されている．
シングルクオートには関数が続き，たとえその関数が定義済みであったとしても，その適用を評価せずシンボリックな関数適用として返す．
\begin{lstlisting}[language=egison,numbers=none]
def f (x: MathValue) : MathValue := x + 1

f a -- a + 1
'f a -- 'f a
\end{lstlisting}
シンボリックな関数適用は，出力でも先頭にシングルクオートを付けて表示される．
シングルクオートは，\ref{casdetail-mathfunc}節で紹介した\lstinline{declare apply}文の本体で，これ以上簡約できない入力に対してシンボリックな値を返すために，\lstinline{sqrt}や三角関数\lstinline{sin}・\lstinline{cos}などの定義のなかで使われている．

\section{\texttt{withSymbols}式によるローカルシンボルの宣言}\label{mexpr-withsym}

\lstinline{declare symbol}文で宣言されるシンボルは大域的である．
一時的にシンボルを使いたいだけの場合のために，ローカルシンボルを宣言するための構文として\lstinline{withSymbols}式がEgisonには用意されている．
\lstinline{withSymbols}式は，第一引数に変数のリストをとり，これらの変数は第二引数の式の評価時にシンボルとして扱われる．
たとえその変数が定義済みであったとしても，\lstinline{withSymbols}式の内部ではシンボルとして扱われる．
\begin{lstlisting}[language=egison,numbers=none]
def k : Integer := 10

k -- 10
withSymbols [k] k -- k
k + withSymbols [k] k -- k + 10
\end{lstlisting}

\lstinline{withSymbols}式で宣言されたローカルなシンボルは，\lstinline{declare symbol}文で宣言されたグローバルなシンボルとは独立したシンボルとして処理される．
以下の2つの\lstinline{j}は同じ名前で表示されるが，別のシンボルであるため，\lstinline{2 * j}のようにまとめられない．
\begin{lstlisting}[language=egison,numbers=none]
declare symbol j

j + withSymbols [j] j -- j + j
\end{lstlisting}

\section{数式データに対するパターンマッチ}\label{mexpr-pattern}

本節は，数式データに対して用意されているパターンを紹介する．
数式に対するマッチャー\lstinline{mathValue}はEgisonライブラリに実装されている．
\lstinline{mathValue}の定義は\texttt{lib/math/expression.egi}で確認できる．
また数式に対するパターンを，より数式に近いかたちで記述するためにいくつかの糖衣構文が組み込みで実装されている．
本節では，これらの糖衣構文も紹介する．
数式に対するパターンマッチは，\ref{cas-derivative}節で実演されている．
パターンの具体的な使用例として，\ref{cas-derivative}節のプログラムを参照してほしい．

\subsection{因子に対するパターンマッチ}

因子に対するパターンコンストラクタには，\lstinline{symbol}・\lstinline{apply1}から\lstinline{apply4}・\lstinline{quote}・\lstinline{func}がある．
\lstinline{symbol}はシンボルに，\lstinline{apply1}から\lstinline{apply4}は\ref{cas-symbol2}節で紹介したシンボリックな関数適用に，\lstinline{quote}は\ref{mexpr-backquote}節で紹介したバッククオートされた式に，\lstinline{func}は第\ref{funcsym}章で解説する関数シンボルと呼ばれるオブジェクトにパターンマッチするためのパターンコンストラクタである．

\lstinline{symbol}パターンコンストラクタの第一引数はシンボルの名前の文字列に，第二引数は添字のリストにパターンマッチする．
添字については，テンソルについて解説する第\ref{tensor}章で改めて紹介する．
\begin{lstlisting}[language=egison,numbers=none]
match x as mathValue with
  | symbol $s _ -> s
-- "x"
\end{lstlisting}
シンボルに対しては\lstinline{symbol}パターンコンストラクタを使ってパターンマッチすることもできるが，実際には値パターンを使ってパターンマッチすることが多い．
\begin{lstlisting}[language=egison,numbers=none]
match x as mathValue with
  | #x -> "Matched"
  | _ -> "Not matched"
-- "Matched"
\end{lstlisting}

シンボリックな関数適用に対しては，引数の個数に応じて\lstinline{apply1}から\lstinline{apply4}のパターンコンストラクタを使う．
第一引数は関数自身に，残りの引数はそれぞれの引数の数式にパターンマッチする．
関数自身に対しては，\lstinline{#sqrt}のような値パターンがよく使われる．
\begin{lstlisting}[language=egison,numbers=none]
match (sqrt x) as mathValue with
  | apply1 #sqrt $g -> g
-- x
\end{lstlisting}

\lstinline{quote}パターンコンストラクタは引数に数式に対するパターンをとる．
このパターンは\lstinline{mathValue}マッチャーを使ってクオートの中身の式とパターンマッチされる．
\begin{lstlisting}[language=egison,numbers=none,escapechar=']
match `(x + 1) as mathValue with
  | quote $e -> e
-- x + 1
\end{lstlisting}


\subsection{項に対するパターンマッチ}

項に対するパターンマッチには，\lstinline{term}パターンコンストラクタを使う．
\lstinline{term}パターンコンストラクタの第一引数は係数に，第二引数は因子と冪指数の連想リストにパターンマッチする．
因子の連想リストのパターンマッチには，\ref{matcher-assoc}節で紹介したassocMultisetマッチャーが使われる．
\begin{lstlisting}[language=egison,numbers=none]
matchAll 3 * x^2 * y as mathValue with
  | term $a (($x, $n) :: $xs) -> (a, x, n, xs)
-- [(3, x, 2, [(y, 1)]), (3, y, 1, [(x, 2)])]
\end{lstlisting}

\lstinline{term}パターンコンストラクタには，より数式に近いパターンを記述するための糖衣構文がある．
以下の1つ目のパターンの\lstinline{$x}は，\lstinline{x^2}のように冪を含めた因子にマッチし，\lstinline{$xs}は残りの因子の積にマッチする．
2つ目のパターンのように，冪指数に対するパターンを\lstinline{^}に続けて書くこともできる．
\begin{lstlisting}[language=egison,numbers=none]
matchAll 3 * x^2 * y as mathValue with
  | $a * $x * $xs -> (a, x, xs)
-- [(3, x^2, y), (3, y, x^2)]

matchAll 3 * x^2 * y as mathValue with
  | $a * $x^#2 * $xs -> (a, x, xs)
-- [(3, x, y)]
\end{lstlisting}

\subsection{多項式に対するパターンマッチ}

多項式に対するパターンマッチには，\lstinline{poly}パターンコンストラクタを使う．
\lstinline{poly}パターンコンストラクタは引数に項のリストにマッチするパターンをとる．
\begin{lstlisting}[language=egison,numbers=none]
matchAll x + y + 1 as mathValue with
  | poly ($x :: $xs) -> (x, xs)
-- [(y, [x, 1]), (x, [y, 1]), (1, [y, x])]
\end{lstlisting}

\lstinline{poly}パターンコンストラクタにも，より数式に近いパターンを記述するための糖衣構文がある．
\begin{lstlisting}[language=egison,numbers=none]
matchAll x + y + 1 as mathValue with
  | $x + $xs -> (x, xs)
-- [(y, x + 1), (x, y + 1), (1, y + x)]
\end{lstlisting}

\subsection{有理式に対するパターンマッチ}

有理式に対するパターンマッチには，\lstinline{frac}パターンコンストラクタを使う．
\lstinline{frac}パターンコンストラクタの第一引数は分子の多項式に，第二引数は分母の多項式にパターンマッチする．
\begin{lstlisting}[language=egison,numbers=none]
matchAll (a + b) / (c + 1) as mathValue with
  | frac $x $y -> (x, y)
-- [(b + a, c + 1)]
\end{lstlisting}
なお，\ref{mexpr-expr}節で述べたように，分母がひとつの項だけからなる分数は負の冪をもつ項として表現されるため，\lstinline{frac}パターンにはマッチしない．

\lstinline{frac}パターンコンストラクタにも，より数式に近いパターンを記述するための糖衣構文がある．
\begin{lstlisting}[language=egison,numbers=none]
matchAll (a + b) / (c + 1) as mathValue with
  | $x / $y -> (x, y)
-- [(b + a, c + 1)]
\end{lstlisting}
